\begin{exo}
	Une entreprise de logistique évalue la performance de ses robots trieurs de colis.

	On étudie un lot de colis à l'aide de deux critères :

	\begin{itemize}
		\item $A$ : \og{} Le colis est fragile \fg{} ;
		\item $B$ : \og{} Le colis est correctement orienté par le robot \fg{}.
	\end{itemize}

	L'ingénieur en chef a mesuré les effectifs mais une partie des données dépend du nombre de robots de type \og{} Alpha \fg{} en activité, noté $n$ (où $n$ est un entier positif).

	\columnbreak

	Le tableau croisé des effectifs est le suivant :

	\begin{center}
		\renewcommand{\arraystretch}{1.2}
		\begin{tabular}{|C{1.5cm}|C{1cm}|C{1cm}|C{1.5cm}|}
			\hline
			               & $B$      & $\overline{B}$ & \textbf{Total} \\ \hline
			$A$            & $n$      & $12$           & $n + 12$       \\ \hline
			$\overline{A}$ & $48$     & $60$           & $108$          \\ \hline
			\textbf{Total} & $n + 48$ & $72$           & $n + 120$      \\ \hline
		\end{tabular}
	\end{center}

	On choisit au hasard un colis dans ce lot.

	\begin{enumerate}
		\item Exprimer en fonction de $n$ les probabilités :
		      \[
			      P(A \cap B)\ \text{et}\ P(B)
		      \]

		\item Rappeler la formule du cours exprimant la probabilité conditionnelle $P_{B}(A)$ en fonction de $P(A \cap B)$ et $P(B)$.
		\item En déduire que la probabilité que le colis soit fragile sachant qu'il est correctement orienté s'exprime par :
		      \[ P_{B}(A) = \frac{n}{n + 48} \]
	\end{enumerate}

	L'ingénieur souhaite que la probabilité qu'un colis correctement orienté soit fragile soit exactement égale à $0,2$ (soit $\frac{1}{5}$).

	\ber
	\item Déterminer la valeur de $n$ pour atteindre cet objectif.
	\end{enumerate}
\end{exo}
